\documentclass[11pt,a4paper]{article}

\usepackage[margin=1.1in]{geometry}
\usepackage{amsmath,amssymb,amsfonts}
\usepackage{algorithm}
\usepackage{algpseudocode}
\usepackage{booktabs}
\usepackage{enumitem}
\usepackage{xcolor}
\usepackage{url}
\usepackage[hidelinks]{hyperref}

\newcommand{\defn}[1]{\textbf{\textit{#1}}}
\newenvironment{keypoint}
  {\begin{quote}\small\noindent\textbf{In plain words.}\ }
  {\end{quote}}

\title{\bfseries Coordination Before Trajectories\\[0.3em]
\large A Ground-Up Explanation of the Problem and the Method}
\author{Working note --- not a submission draft}
\date{}

\begin{document}
\maketitle

\begin{abstract}
\noindent
This is a teaching document, not a paper. It explains the problem we are solving and the
method we built, from the ground up, defining every term the first time it is used. There
are no experiments here: the goal is that a reader who has never seen a motion-planning
paper can follow the whole construction and say precisely where they disagree. Worked
numerical examples use the real robot parameters and have been checked against the
implementation, so every number below is one the code actually produces.
\end{abstract}

\tableofcontents

\bigskip\hrule\bigskip

\part*{Part I --- What Kind of Problem This Is}
\addcontentsline{toc}{part}{Part I --- What Kind of Problem This Is}

\section{Paths, trajectories, and why the difference matters}

A \defn{path} is a curve through space. It says \emph{where} a robot goes and nothing about
\emph{when}. ``Drive from the door to the window along this line'' is a path.

A \defn{trajectory} is a path with timing attached. It says where the robot is at every
instant: at $t = 3.2$\,s it is at $(4.1, 2.7)$ facing north-east and moving at
$0.3$\,m/s. ``Leave at noon, arrive in eight seconds, reaching top speed halfway'' is a
trajectory.

The distinction is the whole difficulty of this subject. Two robots whose \emph{paths} cross
do not necessarily collide --- if one passes the crossing point before the other arrives,
nothing happens. Two robots whose paths never cross can still collide, if their bodies are
wide and the paths run close together. Collision is a property of trajectories, not paths.

\begin{keypoint}
Path = the line on the floor. Trajectory = the line plus a timetable. Robots crash because
of timetables, not lines.
\end{keypoint}

\section{Kinematic versus kinodynamic}

A \defn{kinematic} model of a robot describes what shapes it can move along, ignoring
inertia. A common kinematic model says: this robot can drive forwards and turn, but cannot
slide sideways. Under such a model you may assume the robot can start, stop, and change
direction instantly.

A \defn{kinodynamic} model adds the dynamics: mass, momentum, and bounded actuation. Now
the robot has a \emph{velocity} that is part of its state, it cannot change that velocity
faster than its motors allow, and it therefore cannot stop instantly or turn on the spot
while moving.

The word ``kinodynamic'' in this work means specifically that:
\begin{itemize}[nosep]
\item velocity is part of the robot's state, not something we choose freely at each moment;
\item we choose \emph{accelerations}, and velocity is the accumulated result;
\item accelerations and velocities are both bounded.
\end{itemize}

\begin{keypoint}
Kinematic: ``the robot is here now, put it there next.'' Kinodynamic: ``the robot is here
moving at this speed; I can push the pedal this hard, and it will take a while to stop.''
\end{keypoint}

Why it matters for coordination: a classical multi-robot path-finding method may resolve a
conflict by telling a robot to \emph{wait one step}. A kinodynamic robot moving at
$0.5$\,m/s cannot wait. It must decelerate, which takes time and, crucially, covers
distance. With our parameters that stopping distance is
\[
d_{\mathrm{brake}} \;=\; \frac{v_{\max}^2}{2 a_{\max}}
\;=\; \frac{0.5^2}{2 \times 0.25} \;=\; 0.50\ \text{m},
\]
which is twice the radius of the goal region the robot is trying to stop inside. This robot
is firmly in the regime where ``just stop'' is not an available move.

\section{Describing one robot}

\subsection{State}

The \defn{state} of a system is the smallest collection of numbers that, together with the
future control inputs, determines the future completely. For our robot the state is five
numbers:
\[
\mathbf{x} \;=\; \begin{bmatrix} x & y & \theta & v & \omega \end{bmatrix}^{\!\top}
\]
where $(x,y)$ is the position of the body centre in metres, $\theta$ is the
\defn{heading} (the direction the robot faces, in radians, measured anticlockwise from the
$+x$ axis), $v$ is the forward speed in m/s, and $\omega$ is the \defn{yaw rate} --- how
fast the heading is changing, in rad/s.

Position and heading alone would be a kinematic description. Including $v$ and $\omega$ is
what makes it kinodynamic: the robot ``remembers'' how fast it is going.

\subsection{Control}

The \defn{control} is what we get to choose at each instant:
\[
\mathbf{u} \;=\; \begin{bmatrix} a & \alpha \end{bmatrix}^{\!\top}
\]
where $a$ is \defn{linear acceleration} (how hard we push the speed up or down, in
m/s$^2$) and $\alpha$ is \defn{angular acceleration} (how hard we push the yaw rate up or
down, in rad/s$^2$). We do \emph{not} choose position or speed directly. We choose the two
accelerations, and everything else follows.

\subsection{The equations of motion}

The model is the \defn{second-order unicycle}. ``Unicycle'' because it rolls forward along
its own heading and cannot slide sideways; ``second-order'' because we command
accelerations (second derivatives of position) rather than velocities.

\begin{equation}
\underbrace{
\begin{bmatrix}\dot{x}\\ \dot{y}\\ \dot{\theta}\\ \dot{v}\\ \dot{\omega}\end{bmatrix}
}_{\text{rate of change of state}}
=
\begin{bmatrix}
v\cos\theta \\ v\sin\theta \\ \omega \\ a \\ \alpha
\end{bmatrix}
\label{eq:dyn}
\end{equation}

Reading it line by line. A dot over a symbol means ``rate of change of'':
\begin{itemize}[nosep]
\item $\dot{x} = v\cos\theta$ and $\dot{y} = v\sin\theta$: the robot moves at speed $v$ in
      the direction it is facing. If it faces east ($\theta = 0$) then $\cos\theta = 1$,
      $\sin\theta = 0$, and all the motion is in $x$. \emph{This pair of equations is the
      no-sideways-sliding constraint};
\item $\dot{\theta} = \omega$: the heading changes at the yaw rate, by definition;
\item $\dot{v} = a$ and $\dot{\omega} = \alpha$: our two controls are exactly the rates of
      change of the two velocities.
\end{itemize}

Note what \eqref{eq:dyn} does \emph{not} say: it does not forbid $a$ and $\alpha$ from being
non-zero at the same time. A robot may accelerate and turn simultaneously. We will return
to this, because our method chooses not to, and that choice has a measurable cost.

\subsection{Bounds}

Real motors saturate. Four bounds apply at every instant:
\[
|v| \le v_{\max}, \qquad |\omega| \le \omega_{\max}, \qquad
|a| \le a_{\max}, \qquad |\alpha| \le \alpha_{\max}.
\]
Our values, taken from a published benchmark model so that results are comparable with the
literature, are
\[
v_{\max} = 0.5\ \mathrm{m/s},\quad
\omega_{\max} = 0.5\ \mathrm{rad/s},\quad
a_{\max} = 0.25\ \mathrm{m/s^2},\quad
\alpha_{\max} = 0.25\ \mathrm{rad/s^2}.
\]

Two consequences are worth committing to memory, because they explain most of the method's
awkwardness.

\paragraph{Braking distance.} $d_{\mathrm{brake}} = v_{\max}^2/(2a_{\max}) = 0.50$\,m, as
computed above.

\paragraph{Minimum turning radius.} If a robot moves at speed $v$ while turning at yaw rate
$\omega$, it sweeps a circle of radius $\rho = v/\omega$. To turn tightly you must go
slowly. At full speed the tightest possible circle is
\[
\rho_{\min} \;=\; \frac{v_{\max}}{\omega_{\max}} \;=\; \frac{0.5}{0.5} \;=\; 1.00\ \text{m}.
\]
This number will matter enormously later: the corridors this robot must coordinate inside
are narrower than $1$\,m, so \emph{it cannot corner at full speed inside its own lane}.
Turning radius and speed are not independent, and that coupling is a property of the robot,
not of our algorithm.

\subsection{Time steps and integration}

Equation \eqref{eq:dyn} is continuous. A computer advances it in small steps of
$dt = 0.1$\,s. Taking a continuous equation and stepping it forward numerically is called
\defn{integration}, and the recipe used is the \defn{integrator}.

The simplest integrator, \defn{semi-implicit Euler}, updates velocity first and then uses
the \emph{new} velocity to update position:
\begin{align}
v_{k+1} &= v_k + a_k\, dt, \label{eq:euler-v}\\
x_{k+1} &= x_k + v_{k+1} \cos\theta_{k+1}\, dt. \label{eq:euler-x}
\end{align}
We will use this form for hand calculations because it is exact and easy to reason about.
The simulator itself uses a more accurate recipe (\defn{RK4}, a four-stage
Runge--Kutta method that samples the rate of change at four points within the step and
averages them). The distinction matters only for numerical fidelity, not for understanding.

\begin{keypoint}
A \emph{step} is $0.1$ seconds. A trajectory of ``633 steps'' lasts $63.3$ seconds. When
we say a plan must fit in $1300$ steps, we mean $130$ seconds.
\end{keypoint}

\section{Describing the body and what counts as a crash}

\subsection{The body}

The robot is not a point. It is a \defn{rigid body} --- a shape that moves without
deforming --- and here it is a rectangle of
\[
\text{forward extent } w = 0.5\ \text{m}, \qquad
\text{lateral extent } \ell = 0.25\ \text{m},
\]
meaning $0.5$\,m long in the direction it faces and $0.25$\,m wide across.

\begin{keypoint}
Beware the naming: in the configuration file the field called \texttt{width} is the
\emph{forward} extent and the field called \texttt{length} is the \emph{lateral} extent.
This is inherited from the benchmark and is the reverse of most people's intuition.
\end{keypoint}

\subsection{Two derived numbers}

\paragraph{Bounding radius.} The \defn{bounding radius} is the radius of the smallest circle
centred on the body that contains it --- half the diagonal:
\[
r_{\mathrm{bound}} \;=\; \tfrac{1}{2}\sqrt{w^2 + \ell^2}
\;=\; \tfrac{1}{2}\sqrt{0.5^2 + 0.25^2} \;=\; 0.2795\ \text{m}.
\]
Its use is as a cheap conservative filter. If two body centres are more than
$2 r_{\mathrm{bound}}$ apart, the bodies certainly do not touch, and we can skip the
expensive exact test. If they are closer, we must do the exact test --- they \emph{might}
touch.

\paragraph{Body diagonal.} The \defn{body diagonal} is
\[
d_{\mathrm{diag}} \;=\; \sqrt{w^2 + \ell^2} \;=\; \sqrt{0.5^2 + 0.25^2} \;=\; 0.5590\ \text{m},
\]
which is the largest centre-to-centre distance at which two of these bodies can still touch,
achieved when both are oriented along the line joining their centres. This is the number to
use when asking ``could these two ever conflict, whatever their headings?''

Using the lateral extent $\ell = 0.25$ instead would be a serious error: it is the right
number only for two robots passing side by side, and would declare safe a pair that is
about to collide corner-to-corner.

\subsection{Overlap versus clearance}

Two bodies \defn{overlap} (collide) if they share any point. The
\defn{surface-to-surface distance} between two non-overlapping bodies is the shortest
distance between any point of one and any point of the other; it is $0$ when they touch.

We do not merely require no overlap. We require a \defn{clearance}: the
surface-to-surface distance must stay above $\varepsilon = 0.05$\,m at all times. This is
strictly stronger than non-overlap, and it is the standard the method must meet. A plan that
brings two robots within $0.0122$\,m of each other is rejected even though nothing
technically collides.

\begin{keypoint}
``No overlap'' means the boxes do not intersect. ``Clearance $0.05$'' means there is always
at least a $5$\,cm gap of air between them. We demand the second.
\end{keypoint}

\section{The team problem}

We have $n$ robots (in our scenarios $n \in \{16, 32\}$) in a square workspace $17 \times
17$\,m containing fixed obstacles. Robot $i$ starts at a given state and must reach a goal
position $\mathbf{g}_i$.

Our scenarios are all \defn{permutation swaps}: every robot's goal is some other robot's
start. This is not incidental --- it makes the problem hard in a specific way. In a swap the
place robot $A$ is trying to reach is currently occupied by robot $B$, so $B$ must leave
before $A$ arrives. Coordination is not optional; it is the entire problem.

A \defn{solution} is a set of control sequences $\mathbf{u}_i(\cdot)$ such that
\begin{enumerate}[nosep]
\item every robot ends within the \defn{goal radius} $r_{\mathrm{goal}} = 0.2$\,m of its
      goal;
\item no robot ever overlaps an obstacle;
\item no two robots are ever closer than the clearance $\varepsilon = 0.05$\,m;
\item the whole thing finishes within the \defn{horizon} $T_{\max} = 1300$ steps ($130$\,s).
\end{enumerate}

The \defn{makespan} is the time at which the last robot arrives --- the length of the whole
team's plan. A plan that is collision-free but takes $2000$ steps is not a solution here,
because the episode ends at $1300$.

\bigskip\hrule\bigskip

\part*{Part II --- The Geometric Vocabulary}
\addcontentsline{toc}{part}{Part II --- The Geometric Vocabulary}

This part defines the geometry the method reasons with. None of it is deep, but all of it is
used constantly, and the terms are the ones a reader unfamiliar with the area is most likely
to bounce off.

\section{Routes, polylines, and arclength}

\subsection{Polyline}

A \defn{polyline} is a curve made of straight segments joined end to end, specified by its
list of corner points (\defn{vertices}). Every route in this work is a polyline. It is a
\emph{path}, not a trajectory: no timing yet.

\subsection{Arclength}

\defn{Arclength} is distance measured \emph{along the curve}, as a piece of string laid on
it would measure --- not straight-line distance between endpoints.

For a polyline with vertices $\mathbf{p}_0, \dots, \mathbf{p}_M$, the arclength from the
start to vertex $k$ is the sum of the segment lengths up to that point:
\[
L_k \;=\; \sum_{j=1}^{k} \|\mathbf{p}_j - \mathbf{p}_{j-1}\|,
\qquad L_{\mathrm{total}} = L_M .
\]

\paragraph{Worked example.} Take the polyline $(0,0) \to (3,0) \to (3,4)$. The first
segment has length $3$, the second has length $4$, so $L_{\mathrm{total}} = 7$. The
straight-line distance from start to end is only $\sqrt{3^2+4^2} = 5$. Arclength is $7$; the
robot really does drive $7$\,m.

\subsection{Sampling by arclength, and why not by index}

To \defn{sample} a route means to pick a finite set of points along it. We sample each route
at $m = 64$ points spaced \emph{equally by arclength}: consecutive samples are
$L_{\mathrm{total}}/63$ apart as measured along the curve.

The alternative --- taking the vertices themselves, or spacing samples equally in the
vertex index --- is wrong for our purpose, and the example above shows why. That polyline
has three vertices. Its middle \emph{vertex} is $(3,0)$, sitting at arclength $3$ out of
$7$, which is $43\%$ of the way along. But the true halfway point by distance is at
arclength $3.5$, which is $(3, 0.5)$ --- half a metre up the second segment. If one route
is described with many vertices and another with few, index-based comparison compares
points that are nowhere near equally far along. Arclength sampling removes that
distortion.

\subsection{Normalised progress}

Having sampled by arclength, we label sample $k$ with the \defn{normalised progress}
\[
s_k \;=\; \frac{k}{m-1} \;\in\; [0, 1],
\]
so $s = 0$ is the start of the route, $s = 1$ is the goal, and $s = 0.5$ is the halfway
point \emph{by distance travelled}. ``Normalised'' here just means rescaled to run from $0$
to $1$ regardless of how long the route actually is.

This lets us compare robots with routes of different lengths on a common footing: at the
same $s$, each robot is the same \emph{fraction} of the way through its own journey.

\subsection{What normalised progress assumes --- read this carefully}

The whole point of the common $s$-grid is that we can ask ``are robots $i$ and $j$ close at
the same $s$?'' and treat the answer as a conflict.

That question is about \emph{progress}, not about \emph{time}. Two robots at the same $s$
are equally far along their own routes, but they are not necessarily there at the same
moment. Treating same-$s$ proximity as a conflict silently assumes that all robots advance
along their routes at roughly comparable rates. Where that holds, the conflict set computed
this way is a good stand-in for the true set of space-time conflicts. Where it fails --- a
robot that must turn around first and loses $200$ steps to the turn, for instance --- the
proxy both misses real conflicts and invents false ones.

Two things keep this honest. First, the sampled route is never executed; it exists only to
decide \emph{who contends with whom, and where}. The trajectories are continuous-time
integrations of \eqref{eq:dyn}. Second, the finished plan is checked by simulating it, so a
bad conflict set can produce a bad \emph{assignment}, which verification then rejects --- it
cannot produce an unsafe accepted plan.

\begin{keypoint}
The progress grid is the one genuinely discrete, genuinely approximate step in an otherwise
continuous pipeline. It is a conflict-detection heuristic, not the trajectory
representation. Replacing it with a continuous space-time conflict test is the single most
valuable improvement available.
\end{keypoint}

\section{Tangents and normals}

The \defn{tangent} at a point on a curve is the unit vector pointing along the direction of
travel. On a polyline segment from $\mathbf{p}_k$ to $\mathbf{p}_{k+1}$ it is
\[
\mathbf{t} \;=\; \frac{\mathbf{p}_{k+1} - \mathbf{p}_k}{\|\mathbf{p}_{k+1} - \mathbf{p}_k\|}.
\]
Dividing by the length makes it a \defn{unit vector} (length exactly $1$), so it carries
direction only.

A \defn{normal} is a unit vector perpendicular to the tangent. In two dimensions there are
exactly two, pointing to either side. Writing $\mathbf{t} = (t_x, t_y)$:
\[
\mathbf{n}_{\mathrm{left}} = (-t_y,\ t_x),
\qquad
\mathbf{n}_{\mathrm{right}} = (t_y,\ -t_x).
\]

\paragraph{Worked example.} A robot heading east has $\mathbf{t} = (1,0)$. Its right-hand
normal is $(0,-1)$, which points south. Displacing it ``to its right'' moves it south ---
correct, since south is on the right of someone facing east. A robot heading north has
$\mathbf{t} = (0,1)$ and right-hand normal $(1,0)$, pointing east. Both robots, pushed to
their own right, end up circulating the same way around a point between them. This
observation is the entire roundabout device of Section~\ref{sec:orbit}.

\paragraph{World frame versus body frame.} A \defn{world-frame} direction is fixed relative
to the map (``north''). A \defn{body-frame} direction is relative to the robot (``the
robot's left''). Two robots facing opposite ways have opposite body-frame lefts but the same
world-frame north. Getting this wrong is a real and easy bug: displacing two head-on robots
each ``to its own left'' pushes them both to the \emph{same} side of the road, into each
other.

\section{Graphs}

A \defn{graph} is a set of \defn{nodes} with \defn{edges} joining some pairs. Here nodes are
robots and an edge means ``these two contend for space''; the result is the
\defn{conflict graph}.

A \defn{connected component} is a maximal group of nodes reachable from one another by
following edges. Robots in different components never interact, so each component is
coordinated independently. Note that components grow by \emph{transitivity}: if $A$
conflicts with $B$ and $B$ with $C$, then $A$, $B$, $C$ are one component even if $A$ and
$C$ never come near each other. This will matter.

A \defn{graph colouring} assigns a label (``colour'') to each node such that no edge joins
two nodes of the same colour. We use \defn{greedy colouring}: process nodes in order of
decreasing \defn{degree} (number of edges), giving each the lowest-numbered colour not
already used by a neighbour. It is fast and gives no guarantee of using the fewest possible
colours --- adequate here, because a colour is a lane and we merely need conflicting robots
in different lanes.

\bigskip\hrule\bigskip

\part*{Part III --- The Method}
\addcontentsline{toc}{part}{Part III --- The Method}

\section{The idea, and why it is different}

The standard approach to multi-robot planning with dynamics is: plan each robot separately;
check for conflicts; where two trajectories conflict, add a constraint and \emph{ask a
planner for a new trajectory}. The planner might be a trajectory optimiser (a numerical
solver that adjusts a whole trajectory to minimise a cost subject to constraints), a
sampling-based planner (which grows a random tree of reachable states), or a joint solve
over several robots at once. Methods differ in which planner and which constraints, but the
unit of work is always ``find me a trajectory''.

That question is expensive, and the number of times it is asked grows with the number of
conflicts, which grows with density. In our own reimplementation of a state-of-the-art
ladder of such methods, a $32$-robot instance consumed over $300$ solver calls and a
$1800$-second budget and still returned nothing.

Our claim is that at high density the cheaper question is \emph{``what is the
coordination?''} --- and that the answer is already visible in the shape of the conflict
graph.

\paragraph{The traffic analogy.} A traffic engineer facing a two-way street does not compute
a trajectory per car. They paint a line and assign a side. Facing a junction where many
roads meet at one point, they do not resolve crossings pairwise; they build a roundabout and
impose one circulation direction. Both are \emph{conventions} fixed before anyone moves, and
each exploits a structural fact: in the first case the conflict is a two-sided encounter, in
the second many routes converge on one place.

Our method reads which of these two situations each conflict component is in, applies the
matching convention, and derives the trajectories analytically.

\section{Overview of the pipeline}

\begin{algorithm}[h]
\caption{Conflict-derived constructive planning}
\begin{algorithmic}[1]
\State Build a \textbf{reference route} per robot; sample it by arclength (Sec.~\ref{sec:refs})
\State Build the \textbf{conflict graph} from route proximity (Sec.~\ref{sec:cgraph})
\ForAll{connected components $C$}
  \If{$C$ is a \textbf{junction}} \Comment{Sec.~\ref{sec:device}}
    \State give every member a right-hand \textbf{orbit} \Comment{Sec.~\ref{sec:orbit}}
  \Else
    \State give every member a \textbf{lane} \Comment{Sec.~\ref{sec:lane}}
  \EndIf
\EndFor
\State Displace and simplify routes (Sec.~\ref{sec:simplify})
\State Turn each route into a \textbf{trajectory} of rest-to-rest legs (Sec.~\ref{sec:legs})
\State Choose a \textbf{departure delay} per robot by insertion (Sec.~\ref{sec:sched})
\State \textbf{Verify} by simulation; commit only if it passes (Sec.~\ref{sec:verify})
\end{algorithmic}
\end{algorithm}

No step calls an optimiser or a sampling planner. Steps 1--8 are geometry, graph algorithms,
and closed-form profiles.

\section{Step 1: Reference routes}
\label{sec:refs}

Each robot gets a \defn{reference route}: a rough polyline from its start to its goal,
sampled at $64$ points by arclength as in Part II.

If an obstacle-free kinematic path is available we use it. Otherwise the reference is simply
the straight segment from start to goal. In an empty workspace these coincide, since the
straight line \emph{is} the best route.

The reference is not a plan and is never driven. Its only job is to answer: which robots
contend, and at what progress along each route does the contention happen?

\section{Step 2: The conflict graph}
\label{sec:cgraph}

Robots $i$ and $j$ get an edge when their references come within
$d_{\mathrm{diag}} + \varepsilon = 0.559 + 0.05 = 0.609$\,m at the same normalised progress:
\[
\min_{k} \ \bigl\| \mathbf{r}_i(s_k) - \mathbf{r}_j(s_k) \bigr\| \;<\; d_{\mathrm{diag}} + \varepsilon .
\]

The threshold is the body diagonal, not the lateral extent, for the reason given in
Part~I: two rectangles can touch with centres a diagonal apart. A pair missed here never
receives any coordination assignment at all, so this test must be conservative.

\section{Step 3: Choosing the device}
\label{sec:device}

Here is the central mechanism. For each connected component $C$ with at least three
members, collect the \defn{closest-approach midpoint} of every conflicting pair inside it:
for pair $(i,j)$, find the progress index $k$ at which the two routes are nearest, and take
the midpoint of the two route points there. Call this set $H(C)$.

Now ask: \emph{do all these meeting points sit in essentially one place?} Formally, with
$\bar{h}$ the mean of $H(C)$,
\begin{equation}
C \text{ is a \defn{junction}} \iff
\max_{h \in H(C)} \|h - \bar{h}\| \;<\; d_{\mathrm{diag}} + \varepsilon .
\label{eq:hub}
\end{equation}

\begin{itemize}
\item \textbf{Spread out} $\Rightarrow$ the component is several distinct crossings that
merely happen to be linked through shared members. Give it \textbf{lanes}.
\item \textbf{All in one place} $\Rightarrow$ the component is a junction: every route
passes through one point. Give it a \textbf{roundabout}.
\end{itemize}

Components of size two are never junctions --- a lane already separates two robots.

\paragraph{Why a lane cannot do a junction's job.} A lane is a \emph{two-sided} device. It
resolves an encounter by putting one robot on each side of it. There are exactly two sides.
If thirty-two routes pass through one point, no assignment of two sides separates them; it
is not a matter of choosing better sides. Before the junction device existed, our scheduler
could seat only $1$ of $32$ robots on such a scenario, and no amount of scheduling could
have fixed it.

\begin{keypoint}
Two robots crossing: paint a line, one each side. Thirty-two routes through one point:
build a roundabout. The test in \eqref{eq:hub} is just ``are all the meeting points on top
of each other?''
\end{keypoint}

\section{Step 4a: The lane device}
\label{sec:lane}

\subsection{Which side: colouring}

Greedy-colour the conflict graph so conflicting robots get different colours. Colour index
$c \in \{0, \dots, L-1\}$ becomes a signed offset about the route centre line:
\[
\Delta_c \;=\; \Bigl(c - \tfrac{L-1}{2}\Bigr)\, p ,
\]
where $p$ is the \defn{lane pitch} (centre-to-centre spacing between adjacent lanes). With
$L=2$ and $p=0.5$ the offsets are $-0.25$ and $+0.25$: one robot a quarter-metre left of
the centre line, the other a quarter-metre right.

\subsection{Which direction is ``sideways'': the lane axis}

An offset needs a direction. We compute one \defn{lane axis} shared by everyone in the
group, in the \emph{world frame}.

Collect the tangents $\mathbf{t}_1, \dots, \mathbf{t}_K$ of the group's members at their
contested stretches, and form the $2\times2$ matrix
\[
\mathbf{M} \;=\; \sum_{k} \mathbf{t}_k \mathbf{t}_k^{\!\top}.
\]
This is a \defn{covariance-like scatter matrix}: it summarises which direction the tangents
mostly point along. Its \defn{eigenvectors} are the two perpendicular directions in which
this scatter is purely stretched, and the \defn{eigenvalue} attached to each says how much
scatter lies along it. Taking the eigenvector with the \emph{largest} eigenvalue gives the
dominant direction of travel $\mathbf{d}$; the lane axis is its normal $(-d_y, d_x)$.

\begin{keypoint}
``Eigendecomposition of the tangent scatter'' means: look at all the directions these robots
are travelling, and find the single line that best represents them. Sideways is
perpendicular to that line.
\end{keypoint}

Two robots meeting head-on have tangents differing by $180^\circ$. The matrix
$\mathbf{t}\mathbf{t}^{\!\top}$ is unchanged if $\mathbf{t}$ is negated, so both contribute
the same thing and the axis comes out along the road --- which is what we want. Using each
robot's own body-frame normal instead would flip with heading and push both robots to the
same side; we made exactly that mistake, and it produced $944$ collisions.

\subsection{Splitting a component by direction}

One axis per connected component is right when the component \emph{is} one encounter. It
stops being right when the component has grown by transitivity into a group of robots
travelling in unrelated directions --- then the ``dominant direction'' is an average of
things that should not be averaged, and displacing everyone along its normal separates
nobody.

So we split a component into groups of similar heading before computing an axis. Orientation
is taken \defn{modulo $\pi$} --- that is, directions $180^\circ$ apart are treated as the
same orientation --- so a head-on pair stays together and keeps the shared axis its opposite
lane signs depend on, while genuinely perpendicular encounters are separated.

\subsection{How many lanes actually fit}

The colouring can ask for more lanes than the map holds, and this is a real failure mode
rather than a hypothetical one.

Let $\rho$ be the closest a robot's route comes to a route it does \emph{not} contend with
--- typically the neighbouring corridor. The room available for sideways displacement is
\[
\rho' \;=\; \rho - (\ell + \varepsilon),
\]
subtracting the body width plus clearance that must remain free. We then cap the number of
lanes and compress the pitch only when the nominal width does not fit:
\[
L \;=\; \min\!\Bigl(L_{\mathrm{colour}},\ \Bigl\lfloor \tfrac{\rho'}{\ell + \varepsilon} \Bigr\rfloor + 1 \Bigr),
\qquad
p \;=\; \min\!\Bigl(p_{\mathrm{nominal}},\ \tfrac{\rho'}{L-1}\Bigr).
\]

\paragraph{Worked example.} Rows $1.0$\,m apart, body lateral extent $0.25$, clearance
$0.05$. Then $\rho' = 1.0 - 0.30 = 0.70$\,m of room, and at most
$\lfloor 0.70/0.30 \rfloor + 1 = 3$ lanes fit. If the colouring demands $4$ lanes at the
nominal pitch $0.5$\,m, the outermost offsets would be $\pm 0.75$\,m --- spanning $1.5$\,m
inside a $1.0$\,m corridor. Those robots would be pushed into the neighbouring row, and the
device would \emph{manufacture} conflicts rather than resolve them. Capping to $3$ lanes at
pitch $0.35$\,m keeps the span at $0.70$\,m, exactly the room available.

\subsection{Applying the offset}

The displacement is applied over the stretch of route where the robot is contended, ramped
in and out (\defn{tapered}) so that the start and goal points are untouched. The window is
padded generously: two robots traverse their routes at different rates, so a narrow bulge
can be entered at different times by the two partners, leaving both back on the centre line
when they actually meet. We observed exactly this --- a correct $\pm 0.25$\,m lane
assignment and a dead-centre collision anyway.

If the assigned lane is blocked by an obstacle, the robot \defn{changes lane}: every lane
position is tried, nearest to the assigned one first. The tempting alternative --- widening
the blocked offset --- searches in the wrong direction, since blocked robots are usually in
the outermost lane already and the free space is inboard. Falling back to the centre line is
worse still: that is where the head-on partner is, and it converts a spatial problem into a
timing problem, which is the expensive kind.

\section{Step 4b: The roundabout device}
\label{sec:orbit}

For a junction component with centre $\bar{h}$, \emph{every member is displaced to its own
right}, onto a ring of radius
\begin{equation}
\rho_{\mathrm{ring}} \;=\; \frac{|C| \, (d_{\mathrm{diag}} + \varepsilon)}{2\pi}.
\label{eq:ring}
\end{equation}

\paragraph{Where \eqref{eq:ring} comes from.} In the worst case all $|C|$ robots are on the
ring simultaneously. They sit on a circle of circumference $2\pi\rho$, and each needs
$d_{\mathrm{diag}} + \varepsilon$ of arc to itself. Requiring
$2\pi\rho \ge |C|(d_{\mathrm{diag}} + \varepsilon)$ and solving for $\rho$ gives
\eqref{eq:ring}. It is the smallest ring on which the whole component fits.

\paragraph{Worked example.} With $d_{\mathrm{diag}} + \varepsilon = 0.609$\,m:
\[
|C| = 16 \Rightarrow \rho_{\mathrm{ring}} = \frac{16 \times 0.609}{2\pi} = 1.551\ \text{m},
\qquad
|C| = 32 \Rightarrow \rho_{\mathrm{ring}} = \frac{32 \times 0.609}{2\pi} = 3.102\ \text{m}.
\]

\paragraph{Why ``own right'' and not a shared axis.} This is the exact opposite of the lane
case, and the contrast is the heart of the method. For a lane, a \emph{shared} axis is
essential: it is what makes ``opposite signs'' mean ``opposite sides of the road''. For a
junction, a shared axis is meaningless, because ``the same side'' of a point is a different
direction depending on which way you approach it. As computed in Part~II, a robot heading
east displaced to its right goes south of the centre; a robot heading north displaced to its
right goes east of it. Follow that rule for every robot and the whole component circulates
the same way --- a roundabout emerges, without anyone being told which way to go round.

Because the circulation is collective, mirroring is available to a lane as a fallback but
never to an orbit: sending one member the other way round puts it head-on into the entire
circulation.

\paragraph{A counter-intuitive fact.} One might size the ring generously for safety. Doing
so makes matters \emph{worse}: at scale $1.0$ (the geometric minimum) all $32$ robots can be
scheduled, while at $1.3$ only $21$ can and at $2.0$ only $18$. A wider ring is a longer
detour through the same congested annulus, so the extra room costs more time in the
contested region than the separation buys. The minimum is also the optimum.

\section{Step 5: Simplifying the route}
\label{sec:simplify}

The displaced route has $64$ closely spaced points. Driving it literally would mean $64$
tiny legs. We reduce it to its essential corners with the \defn{Douglas--Peucker}
algorithm: draw a chord from the first point to the last; find the point furthest from that
chord; if that distance exceeds a tolerance, keep the point and recurse on the two halves;
otherwise discard everything between. What survives is the set of points needed to represent
the curve to within the tolerance.

The obvious alternative --- delete any point whose turn angle is below a threshold --- is
unusable here, and instructively so. A lane displacement is \emph{gradual} by construction:
spread over dozens of samples, its per-step turn is tiny everywhere. An angle threshold
therefore deletes the entire lane bulge and puts the robot back on the collision course the
lane was created to avoid. Douglas--Peucker measures deviation from the chord, which is
exactly the quantity a lane is made of, so it preserves it.

\section{Step 6: From route to trajectory}
\label{sec:legs}

Now we attach timing, and this is where the dynamics finally enter.

Each remaining waypoint is reached by one \defn{leg} in two phases:
\begin{enumerate}[nosep]
\item \textbf{turn in place} to face the next waypoint, starting and ending at rest;
\item \textbf{drive straight} to it, starting and ending at rest.
\end{enumerate}
This is called \defn{rest-to-rest}: every phase begins and ends stationary.

\subsection{The trapezoidal profile}

Each phase is one-dimensional: rotate through angle $\Delta\theta$, or translate through
distance $d$, starting and ending at rest, respecting an acceleration bound $\bar{a}$ and a
velocity bound $\bar{v}$.

\paragraph{Triangular case.} Accelerate at $a$ for $n$ steps, then decelerate at $-a$ for
$n$ steps. Under semi-implicit Euler \eqref{eq:euler-v}--\eqref{eq:euler-x}, the distance
covered while accelerating from rest is
\[
\sum_{k=1}^{n} (k\,a\,dt)\,dt \;=\; a\,dt^2\,\frac{n(n+1)}{2},
\]
and while decelerating symmetrically it is $a\,dt^2\,\frac{n(n-1)}{2}$. Adding them, the
cross terms cancel neatly:
\begin{equation}
d \;=\; a\,dt^2\,n^2 .
\label{eq:tri}
\end{equation}
So the smallest feasible $n$ is $\lceil \sqrt{d/(dt^2\bar{a})}\,\rceil$, and $a$ follows from
\eqref{eq:tri}. Peak speed is $n\,a\,dt$; if that exceeds $\bar{v}$ we need the next case.

\paragraph{Trapezoidal case.} Accelerate for $n_u$ steps to exactly $\bar{v}$, \defn{cruise}
at constant speed for $m$ steps, then decelerate for $n_u$ steps. Here
$n_u = \lceil \bar{v}/(\bar{a}\,dt) \rceil$ and the same bookkeeping gives
\begin{equation}
d \;=\; a\,dt^2\,n_u\,(n_u + m).
\label{eq:trap}
\end{equation}
Solving for $m$ gives the cruise length. The name is from the velocity-versus-time plot:
up, flat, down --- a trapezoid. Omitting the cruise phase, as our first implementation did,
forces the whole distance into a triangle and is dramatically slower.

\paragraph{Worked example: driving $15$\,m.} With $dt = 0.1$, $\bar{a} = 0.25$,
$\bar{v} = 0.5$. The triangular solution needs $n = \lceil\sqrt{15/0.0025}\,\rceil = 78$
steps, giving peak speed $78 \times 0.2465 \times 0.1 = 1.92$\,m/s --- far over the
$0.5$ bound, so it is infeasible. The trapezoid uses
$n_u = \lceil 0.5/0.025 \rceil = 20$ steps to reach $0.5$\,m/s, then from \eqref{eq:trap}
$m = 280$ cruise steps, then $20$ to stop: \textbf{$320$ steps} ($32$\,s) at exactly
$a = 0.25$ and peak speed exactly $0.5$. Integrating the resulting control sequence gives
$15.000000$\,m and a final speed of $-2.8\times10^{-17}$, i.e.\ exactly at rest.

\paragraph{Worked example: turning $180^\circ$.} Same arithmetic with
$\Delta\theta = \pi$, $\bar{a} = \alpha_{\max} = 0.25$, $\bar{v} = \omega_{\max} = 0.5$
gives $20 + 43 + 20 = \textbf{83 steps}$ ($8.3$\,s), integrating to exactly $180.0000^\circ$.
A robot that starts facing away from its goal spends $8.3$ seconds just turning round.

\subsection{Why the two phases are separated --- and what it costs}

Splitting each leg into turn-then-drive has two justifications.

First, it makes the closed form possible. During the turn $v = 0$, so position does not
change; during the drive $\omega = 0$, so heading does not change. The two one-dimensional
profiles are therefore completely decoupled, and each is exactly solvable. If both
accelerations fire together, the equations couple --- the heading changes while the robot
translates, so the traversed path curves --- and hitting an exact waypoint becomes a
\defn{two-point boundary value problem}, which is precisely the kind of thing a numerical
optimiser is needed for. Avoiding that call is the entire point of the method.

Second, it guarantees that the path actually driven \emph{is} the polyline on which all the
lane and clearance reasoning was done. If the robot arcs through corners it leaves its lane,
and the clearance argument no longer applies to the thing being executed.

The cost is real and we measure it: across our scenarios, \emph{zero} control steps out of
tens of thousands have both $a \ne 0$ and $\alpha \ne 0$; $21$--$33\%$ of moving steps are
spent turning in place; robots come to a full stop $3$--$6$ times mid-route; and the
resulting horizon is $1.70$--$2.08\times$ a cruise-limited lower bound for the same path.

\subsection{Could both fire at once? Yes --- but the robot fights back}

Nothing in \eqref{eq:dyn} forbids simultaneous $a$ and $\alpha$. The obstacle is the
turning-radius coupling from Part~I: cornering on radius $\rho$ requires
$v \le \omega_{\max}\rho$. Our lanes are $0.35$--$0.50$\,m wide, so cornering \emph{inside a
lane} demands
\[
v \;\le\; 0.5 \times 0.35 \;=\; 0.175\ \text{m/s},
\]
about a third of top speed. A robot simply cannot corner at speed within its own lane; it
must slow down substantially either way.

There is nonetheless clear headroom at the small end. Many turns are not real corners at
all but artefacts of route simplification --- in one scenario the median turn is
$6.7^\circ$, and $62\%$ of turns are under $30^\circ$. Coming to a complete stop to rotate
by seven degrees is plainly wasteful. Steering through small turns while driving, and
keeping stop-and-go for genuine corners, would recover much of the penalty: the lateral
deviation from the chord for a turn $\delta$ over a leg of length $L$ is roughly $L\delta/8$
--- about $0.05$\,m for $7^\circ$ over $5$\,m, comfortably inside a lane, so the clearance
argument survives. This is the clearest improvement available and is not yet implemented.

\section{Step 7: Scheduling departures}
\label{sec:sched}

Lanes are spatial and can be decided from routes alone. \emph{When} each robot sets off
cannot: two robots share a place only if they are there at the same time, and how long a leg
takes is known only once the leg is built. So scheduling happens last, on the finished
trajectories.

Each robot gets a \defn{departure delay} $\delta_i$ --- it sits still at its start for
$\delta_i$ steps, then executes its trajectory.

\subsection{Insertion}

We place robots one at a time. For each, we try delays $0, 5, 10, \dots$ and take the
\emph{smallest} that keeps it clear of every robot already placed. Because we only ever test
against already-committed trajectories, the schedule is correct \emph{by construction}; we
never need to re-detect conflicts and patch them, an approach we tried and abandoned
because applying a shift changes who meets whom, invalidating the very detection that
prompted it.

\paragraph{Occupancy includes standing still.} When testing two robots we compare their
positions over the whole span, treating a robot as sitting at its start before it departs
and at its goal forever after it arrives. Omitting the tail permits a plan in which a robot
parks precisely where a later robot needs to drive; we observed $79$ such collisions before
fixing this.

\subsection{Why greedy insertion gets stuck}

Insertion's only lever is delay --- it can make a robot go \emph{later}, never earlier. In a
permutation swap this is sometimes exactly backwards: robot $A$'s goal is robot $B$'s start,
so $B$ must be \emph{gone before} $A$ arrives. If $B$ is placed late in the order and all
small delays are taken, its only available move --- waiting longer --- makes things worse,
and the run dead-ends with no way back.

Two targeted repairs address this. Both are held back rather than applied first, for a
reason worth stating: each widens the set of feasible schedules, and so changes which
schedule the procedure returns first --- and the rescued schedule is not necessarily the
better one. We observed a repair ``fix'' a scenario by returning a plan using the entire
$1300$-step budget where the unrepaired procedure had been finding one of $1212$ steps. So:
try everything without repairs first, then allow repairs.

\paragraph{Repair 1: precedence from resting places.} Robot $r$ is ordered before $q$ if
$q$'s goal lies on $r$'s route (so $q$, once parked, would block $r$), or if $r$'s start lies
on $q$'s route (so $r$ must clear out before $q$ arrives). Swaps produce two-cycles --- each
partner's goal is the other's start --- which is not a contradiction, since both simply
leave early; cycles are broken by route length, and mutual partners kept adjacent in the
order.

\paragraph{Repair 2: single-blocker eviction.} If a robot cannot be seated at any delay, find
the delay at which it is \emph{least} obstructed. If exactly one already-placed robot stands
in the way there, evict that robot, seat the stranded one, and re-queue the evicted robot.
Each robot may be evicted a bounded number of times, which keeps the procedure finite. This
is one level of backtracking, aimed at a specific measured failure, not a general search.

\subsection{Honest status}

Insertion scheduling is \defn{greedy} --- it commits to each decision without reconsidering
--- and offers no \defn{completeness} guarantee: it may fail on instances that do have
solutions. The repairs narrow the gap; they do not close it.

\section{Step 8: Verification, and the safety contract}
\label{sec:verify}

The assembled plan is checked by the simulator's own collision routines, not by the
planner's internal model: every pair of robots at every shared time index, every robot
against every obstacle, every final position against the true goal radius, and the total
length against the horizon. If anything fails, the construction returns nothing, along with
a machine-readable reason.

This gives a useful contract. Because nothing is committed that has not passed the real
checker, the construction can be installed \emph{ahead} of a conventional solver ladder as a
zeroth rung. If it succeeds, the expensive machinery is never entered. If it fails, the
ladder runs exactly as it would have. It can add solutions; it cannot remove one. A wrong
coordination guess costs one construction and nothing else.

\section{Summary of the argument}

\begin{enumerate}
\item Collision is a property of trajectories, not paths, and a robot with bounded
acceleration cannot stop or turn freely --- so coordination cannot be an afterthought.
\item Conventional methods resolve each conflict by asking a planner for a new trajectory.
That question is expensive and is asked once per conflict, so the cost grows with density.
\item The shape of the conflict graph already says what kind of coordination is needed.
Scattered crossings want lanes; routes converging on a point want a roundabout. The test is
just whether the pairwise meeting points are on top of one another.
\item These two devices are not interchangeable, and their rules are opposites: a lane needs
one \emph{shared} axis for the whole group; a roundabout needs each robot displaced along its
\emph{own} right-hand normal.
\item Once the device is chosen, trajectories follow in closed form --- trapezoidal profiles
under the true bounds --- with no solver.
\item Timing is settled last, by inserting robots one at a time at the least delay that
keeps them clear.
\item Everything is checked by simulation before it is believed, so the construction is safe
to bolt onto an existing planner.
\end{enumerate}

\section*{Known weaknesses, stated plainly}

\begin{itemize}
\item \textbf{The progress grid is discrete and approximate.} Conflicts are found by
comparing robots at equal \emph{progress}, which assumes comparable rates of advance.
Replacing it with a continuous space-time conflict test is the most valuable single
improvement available.
\item \textbf{Controls are interleaved.} Linear and angular acceleration never fire
together, costing $1.70$--$2.08\times$ the cruise-limited bound. The fix for small turns is
identified and unimplemented; large turns are genuinely constrained by the robot's
turning-radius coupling.
\item \textbf{The scheduler is greedy and incomplete.}
\item \textbf{Only two devices exist.} A component whose lane demand exceeds its corridor
capacity has no device that fits, and we have one benchmark scenario in exactly that
position.
\end{itemize}

\end{document}
